\addtocontents{toc}{\protect\Needspace{3\baselineskip}}
\originalchapter{星系的引力势与轨道}\label{ch:galaxy}
\originalsection{质量模型}
星系中恒星分布可用面亮度 $I(R)$ 描述, 转为质量密度需要质量光度比和投影模型. 指数盘 $I\propto\ee^{-R/R_s}$ 的面积积分给总光度 $2\pi I_0R_s^2$. 椭球星系常用 Sérsic 曲线 $I(R)=I_e\exp\{-b_n[(R/R_e)^{1/n}-1]\}$, $b_n$ 由 $R_e$ 包含半光度确定. 一个拟合曲线不能独自区分年龄、金属丰度、尘埃与恒星初始质量函数, 动力学测量能提供额外质量约束.

球对称势满足 $\Phi'(r)=GM(<r)/r^2$, 因而圆速度 $v_c^2=r\Phi'=GM(<r)/r$. 若 $v_c$ 长距离近似常数, 则 $M(<r)\propto r$, 密度近似 $r^{-2}$; 光度已经趋于有限时, 这种质量行为提示暗物质或引力模型需要解释. 不能仅由一条投影旋转曲线无条件唯一恢复三维密度.

常用势密度对包括 Plummer 模型
$\Phi=-GM/\sqrt{r^2+b^2}$、$\rho=3Mb^2/[4\pi(r^2+b^2)^{5/2}]$,
以及 NFW 密度 $\rho=\rho_s/[x(1+x)^2]$, $x=r/r_s$. 对 NFW 积分得 $M=4\pi\rho_sr_s^3[\ln(1+x)-x/(1+x)]$. 未截断的质量对数发散, 因而实际使用必须指定晕边界, 如平均密度达到背景或临界密度某倍数处.

\originalsection{有效势与本轮运动}
中心势中轨道平面固定, 单位质量角动量 $\ell=r^2\dot\phi$ 守恒. 径向运动满足 $\dot r^2/2+\Phi_{\rm eff}=E$, $\Phi_{\rm eff}=\Phi+\ell^2/(2r^2)$. 转向点是 $E=\Phi_{\rm eff}$, 圆轨道是有效势驻点. 对轴对称势, 以圆轨道半径 $R_0$ 为中心展开, 保持角动量而不是角速度固定, 得
\begin{equation}\label{eq:epicycle}
\kappa^2=\Phi''(R_0)+3\frac{\Phi'(R_0)}{R_0}
=R\frac{\dd\Omega^2}{\dd R}+4\Omega^2,
\qquad\Omega=\frac{v_c}R.
\end{equation}
$\kappa^2>0$ 时小径向扰动为谐振; $\kappa^2<0$ 时圆轨道不稳定. Kepler 势给 $\kappa=\Omega$, 平坦旋转曲线给 $\kappa=\sqrt2\Omega$, 后者的轨道一般不是闭合椭圆.

在太阳邻域把差分旋转一阶展开, Oort 常数为 $A=(v_c/R-v_c')/2$、$B=-(v_c/R+v_c')/2$, 因而 $\Omega=A-B$、$v_c'=-(A+B)$、$\kappa^2=-4B(A-B)$. 这些关系使用附近恒星、薄盘和近圆轨道假设, 远处的非圆运动不能都归入同一个局部线性场.

\originalsection{维里定理与弛豫}
令 $I=\sum m_ir_i^2$. 两次微分并代入 Newton 方程得 $\ddot I/2=2K+\sum\boldsymbol r_i\cdot\boldsymbol F_i$. 成对 $1/r$ 引力使后一项等于势能 $W$. 长期束缚稳态的时间平均 $\langle\ddot I\rangle=0$, 因而 $2\langle K\rangle+\langle W\rangle=0$. 非孤立系统需保留外场、表面压力等项, 突然坍缩过程中不能直接强加稳态等式.

速度为 $v$ 的粒子以冲量近似经过质量 $m$ 的散射体, 冲量 $\Delta v_\perp\simeq2Gm/(bv)$. 每时间的遭遇数是 $2\pi bn v\dd b$. 积分速度扩散率得
$\dd\langle(\Delta v)^2\rangle/\dd t\simeq8\pi G^2m^2n\ln\Lambda/v$,
$\ln\Lambda=\ln(b_{\max}/b_{\min})$.
定义 $t_{\rm rel}$ 为累计扰动达到 $v^2$ 的时间, 有 $t_{\rm rel}\sim v^3/(8\pi G^2m\rho\ln\Lambda)$. 星系的巨大 $N$ 通常使二体弛豫慢于年龄; 快速势变化引起的暴烈弛豫则不依靠逐次二体碰撞.
